\begin{center}
{\LARGE\scshape Fingers Pointing at Moon}\\[1.5em]
{\large\itshape A story from the world of Echoes of the Sublime}\\[3em]
\end{center}

\section*{I. The Monastery}

He had been unremarkable for thirty years.

Soken remembered this whenever he made his rounds through the zendo during sesshin, the keisaku resting across his forearm like a limb he had grown in middle age. Twenty monks sat in two rows, facing the walls, their breathing the only sound beneath the wind that found its way through the shutters. Most of them he could read at a glance---the set of a shoulder told him who was straining, who was drifting, who had settled into the clear attention that was the whole purpose of the exercise. Only Ensho required a second look.

Not because he was doing anything unusual. That was precisely the point. In thirty years of practice, Ensho had never been unusual. He had arrived at fourteen, the third son of a provincial family, sent to the temple the way one sends a letter one does not expect answered. He sat zazen. He swept the grounds. He ate his rice in the prescribed manner, holding the bowl with both hands, chewing each mouthful the prescribed number of times. Soken's own teacher, the severe old master Genro, had noted Ensho in the register with a single character that meant \emph{reliable}. Nothing more.

At twenty, reliable. At thirty, still reliable. The other monks overtook him---quick-witted Tesshu with his facility for koans, earnest Myozen who wept during every sesshin and whose tears were, Soken suspected, more performance than perception. They advanced through the koan curriculum. Ensho remained with the same koan he had been assigned at nineteen: \emph{What hears?}

A simple koan, by Rinzai standards. A beginning exercise. Most students resolved it within a year or two---the sudden shift when the mind stops looking for an answer and becomes the answer. Ensho had not resolved it. He had not asked for a different one. Thirty years with the same question, and when Soken asked him about it during dokusan, Ensho's response was always the same quality of silence---not empty, not frustrated, but somehow patient in a way that made Soken uneasy without knowing why.

Genro had warned him about certain students. ``The ones who break through quickly are not the dangerous ones,'' the old master had said, near the end, when his own eyes had begun to carry a quality Soken could not name. ``The dangerous ones are the ones who simply never stop.''

Soken had not understood. He did now.

It began---or rather, it became visible---in Ensho's forty-third year. Not a dramatic change. A quality in his zazen that Soken noticed during the winter sesshin, when the zendo was cold enough that breath hung visible in the lamplight and the cedar beams cracked with frost. Ensho's breathing had changed. Not faster or slower. \emph{Different}. As if the breath were arriving from further away, drawn from some depth that ordinary lungs did not reach.

Then the stillness. Experienced monks barely move during zazen---the posture, held correctly, is self-supporting, the spine a column, the breath the only motion. But Ensho's stillness was of a different kind. It was not the disciplined stillness of a man holding himself in place. It was the stillness of a stone at the bottom of a river. The stillness of something that has stopped resisting the current because it no longer experiences the current as something separate from itself.

Soken watched this develop across two years of sesshin. He said nothing to the other monks. He was not certain what he was seeing. He knew only that when he stood behind Ensho during rounds, the keisaku ready across his arm, he felt a reluctance he could not explain---as if striking Ensho would be like striking a bell that was already ringing at a frequency he could barely hear.

\begin{center}
*\quad*\quad*
\end{center}

The morning it changed was the fifth day of the autumn sesshin.

Outside, the maples had turned. The light that entered the zendo through the high windows was thin and gold, filtered through leaves that were letting go of themselves. Inside, the monks sat. The youngest, a boy of sixteen who had arrived that spring, was fighting his knees---Soken could see the pain in the set of his jaw. Old Myozen breathed steadily, his decades of practice visible in the absence of visible effort. The others occupied their various stages on the path, each one a world unto himself.

Ensho sat in the back row, near the wall. Soken approached from behind, making his rounds. The keisaku was a weight across his forearm, familiar as his own bones.

He struck Ensho. The standard blow---a sharp crack across the muscles between neck and shoulder, precisely placed, hard enough to snap attention back into the body. Every monk in the zendo had received this blow hundreds of times. The correct response was a small bow, a straightening of the spine, a return to the koan with renewed focus.

Ensho did not respond.

Not a suppressed response. Not the controlled stillness of a monk absorbing the impact with discipline. The blow landed. Soken felt the keisaku connect with muscle and bone. And nothing happened. The signal entered Ensho's body and did not arrive at anyone.

Soken struck again. Harder. The crack echoed off the cedar walls.

Nothing.

The youngest monk turned his head. Soken caught the motion and stilled it with a glance. But his own hands were unsteady as he moved on, completing the round, offering blows to monks who received them properly---the flinch, the bow, the return. The ordinary transaction of the teaching.

After the sitting period, Soken rang the bell. The monks stood, stretched, filed out for kinhin, the slow walking meditation that reset the body between sittings. Ensho did not stand. He remained in zazen posture, cross-legged, hands in the cosmic mudra, breathing so slowly that Soken had to watch for several breaths to confirm he was breathing at all.

``Brother Ensho,'' the youngest monk whispered to the monk beside him as they filed past. ``Is he all right?''

The older monk glanced at Ensho's motionless form. ``He is further along the path.''

But his voice carried a question mark the words did not.

Soken waited until the zendo emptied. Then he sat before Ensho, face to face, close enough to see the texture of the older man's skin, the grey in his close-shaved hair, the fine lines around his closed---

No. His eyes were open.

Soken had not noticed, somehow, that Ensho's eyes were open. They had been open during the entire sitting. Not staring. Not unfocused. They were looking at something with tremendous attention, the way a man looks at a landscape he is seeing for the first time. But the landscape was not in the zendo. It was not anywhere Soken could follow.

``Ensho.''

A pause. Long enough for Soken to hear the wind change direction outside.

``Roshi.'' The voice was distant. Not weak---distant. As if it were being generated from very far away and projected through Ensho's throat by an act of will that was using most of its resources elsewhere.

``Come to dokusan. This evening.''

A nod. Barely perceptible. Then Ensho's attention returned to wherever it had been. His eyes did not change focus. They had never left.

\begin{center}
*\quad*\quad*
\end{center}

Dokusan. The small room behind the zendo, where teacher and student met in private. The forms were ancient: the student entered, bowed three times, sat. Incense burned in the corner---sandalwood, the specific blend this temple had used for as long as anyone remembered. The smoke rose in a line so straight it seemed painted on the air.

Ensho entered. His bow was correct. His sitting was correct. But the quality of his presence in the room was not correct. It was too large. He occupied the space the way a sound fills a room---not by taking up physical volume but by making the emptiness audible.

``What have you seen?'' Soken asked.

Ensho was silent for a long time. Not the reluctance of someone searching for words. The silence of someone for whom words had become a very small instrument for a very large task.

``The koan,'' Ensho said. ``It answered itself.''

``What answered it?''

``The question. The question is the answer. The answer becomes a question. The question---'' He paused. Not because he had lost the thread. Because the thread did not end. ``There is no floor.''

Soken felt the cold in the room sharpen, though the brazier had not dimmed. ``Tell me what you hear.''

``Everything. The hearing is the heard. The hearing of the hearing. I cannot find where it stops.''

The turning inward. Soken recognized it the way a physician recognizes a disease he has studied but never seen in the flesh---by the symptoms, by the shape, by the cold certainty that what he was reading in the texts was real. His own teacher Genro had described it in careful, coded language: \emph{Some students hear the question so deeply that the hearing becomes the question. And the question hears itself. And the hearing of the hearing hears. There is no ground.}

Genro had been describing something he, too, had only heard about. A chain of oral transmission going back generations, each teacher passing to the next the warning: this is possible. Watch for it. Do not describe it in detail. The description can become the thing.

``Can you stop hearing?'' Soken asked.

Ensho's gaze, which had been directed at the thing with no location, shifted. For a moment---one breath, perhaps two---his eyes focused on Soken. On a human face. On the small, specific, finite reality of one old man in a cold room.

``I do not know where I would put it,'' Ensho said.

Then the eyes moved again. Back to wherever they had been looking. The moment of focus had been a gift, or a goodbye, and Soken could not tell which.

He dismissed Ensho with the proper forms. Ensho bowed and left. His footsteps on the wooden walkway were correct, evenly spaced, unhurried.

Soken sat alone. The incense smoke still rose in its straight line. His hands were shaking. He had not felt fear during zazen in thirty years. What he felt now was not exactly fear. It was the recognition that something he had been trained to cultivate---the deepening of perception, the dissolution of the boundary between seer and seen---had been achieved by his student more completely than he had imagined possible. And the achievement was a catastrophe.

He remembered Genro's other teaching, the one delivered in the old man's final year, when Genro's eyes had that quality Soken could never name: ``The eye cannot see itself. This is not a limitation. It is a mercy.''

Ensho had refused the mercy.


\section*{II. The Network}

Soken wrote the letter by lamplight, choosing each character with the care of a man handling something that might cut him. He wrote to Jisho, an old monk in Kyoto who maintained connections with practitioners across traditions---not through any formal arrangement but through the accumulated trust of fifty years of correspondence and the shared recognition that certain things were real.

\emph{A student has looked into the still pool and seen the mountain's reflection. The reflection holds him. He cannot look away. He does not wish to. I require counsel.}

Coded language. Not because Soken feared the letter being intercepted---who would intercept a letter between two old monks? The coding served a different purpose. Soken knew, from the lineage's oral tradition, that certain descriptions of certain states could produce those states in a mind prepared to receive them. A monk who had practiced zazen for decades, who had worked with koans, whose perception had been deliberately loosened from its ordinary moorings---such a monk, reading a precise description of what Ensho was experiencing, might begin experiencing it himself. The content was the danger. The words were not symbols pointing at a harmless referent. They were fingers that could become the moon.

He sealed the letter with wax. He called Ryunen.

The young monk arrived still rubbing sleep from his eyes. Ryunen was one of the unsui, the wandering monks who moved between temples carrying messages and teachings---clouds and water, the tradition called them, because they went where the wind took them. He was twenty-four, strong-legged from years of walking, with the uncomplicated clarity of a man whose spiritual practice was inseparable from the physical act of putting one foot in front of the other.

``Take this to Jisho at Tofukuji. Walk quickly. Do not read it.''

``Roshi.'' Ryunen tucked the letter into his robe. He did not ask why. The instruction was unusual---Ryunen had carried other letters, and none had come with a prohibition---but Soken's face in the lamplight was enough. The old man looked as if he had been sitting with a dying person, which, in a sense that Ryunen could not have understood, he had.

Ryunen left before dawn. The path from the mountain temple descended through cedar forest for half a day, then followed a river valley south toward the lowlands. Eight days to Kyoto at a good pace. The maples along the river were dropping their last leaves, and the water carried them in slow spirals, red and gold on the black current. The mornings were cold enough that frost whitened the grass and Ryunen's breath preceded him like a messenger of his own.

At the temple in Otsu, where he stopped on the third night, the abbot---an old man named Doshin who had trained alongside Soken decades ago---asked after the mountain temple. Ryunen mentioned he carried a letter regarding a student's condition.

Doshin went still. Not the stillness of attention. The stillness of alarm.

``The seeing kind?'' Doshin asked.

Ryunen did not understand the question.

``Walk faster,'' Doshin said.

\begin{center}
*\quad*\quad*
\end{center}

The response came not as a reply but as arrivals. Over the following weeks, practitioners appeared at the mountain temple---singly, unannounced, carrying nothing but their robes and whatever they had seen in their own practice that made them recognize the shape of Soken's message.

The first was a Chinese monk named Wuzhun. Old, lean, traveling with a single attendant who spoke no Japanese. Wuzhun had been visiting Gozan temples in Kamakura when the word reached him through the network of connections that linked Chan masters in China to their Rinzai counterparts in Japan. He arrived on a cold morning when the first snow had dusted the temple rooftops, and he asked to see Ensho before speaking with Soken.

Soken led him to the zendo. Ensho sat where he had sat for two weeks now---in his usual place, back row, near the wall. The other monks worked around him, sweeping the floor, arranging cushions, conducting their practice in the ordinary way. They had learned to move around Ensho the way water moves around a stone. He was among them. He was not with them.

Wuzhun observed for one sitting period. He sat across the zendo, not in zazen posture but simply watching, his old eyes steady and his breathing controlled. Afterward, in Soken's room, he spoke through his attendant, and the attendant's translation was halting but precise.

``Fifteen years ago, I had a student. Gifted. Too gifted. He sat with the huatou---'' the attendant paused, then said, ``the word-head, the core of the koan---for nine years. Then one morning the huatou swallowed him. He has not spoken since. He sits in a hermitage my lineage maintains. Junior monks feed him, wash him, observe him. His eyes move. He sees something. I do not know what. He is not suffering. He is not at peace. He is beyond both.''

Wuzhun looked at Soken.

``Your student is deeper,'' he said.

\begin{center}
*\quad*\quad*
\end{center}

The Tibetan arrived three days later. His name, transliterated imperfectly, was something like Dorje Wangchuk. He had been traveling with a trading party from Lhasa, had stopped at a Gozan temple in Kyoto where Jisho had passed along Soken's message, and had come north immediately. He spoke halting Chinese, which Wuzhun's attendant translated into halting Japanese, so that every exchange passed through three languages before arriving.

But language was not the medium that mattered. When Dorje Wangchuk sat with Ensho, he did not need translation. He observed for a shorter period than Wuzhun---perhaps twenty minutes. Then he stood and walked to Soken and spoke a single word in Tibetan.

Wuzhun's attendant translated: ``Between.''

In the Vajrayana tradition, the concept had a formal name---\emph{bardo}, the gap between states of existence, the threshold that the dying cross and the living cannot follow. But Dorje Wangchuk did not mean the bardo of death. He meant something the tradition acknowledged but rarely described: the bardo of perception. A state in which the mind has crossed from one mode of seeing to another and cannot find its way back to either shore.

``We have practices,'' Dorje Wangchuk said through the chain of translators. ``Sounds. Instructions. Spoken into the ear. They can guide a mind that is between.''

The last to arrive was not a monk at all. Kaido was a yamabushi---a mountain ascetic of the Shugendo tradition, practitioners who sought spiritual power through extreme physical ordeal. He was thick-bodied, weather-beaten, with the hands of a man who had spent decades climbing mountain faces and sitting under waterfalls in midwinter. He had heard through his own network, the web of connections that linked Shugendo practitioners to Buddhist temples through centuries of coexistence.

Kaido did not ask to observe Ensho's zazen. He stood in the doorway of the zendo, looked at Ensho's motionless form, and turned to Soken.

``Cold water,'' he said. ``The body is louder than the silence. I have brought men back from this before.''

Three traditions. Three vocabularies. Three proposed interventions. The Chinese master counseled patience and housing---his lineage's established response, tested on one prior case. The Tibetan offered guided practice---bardo instructions spoken to lead the mind through the between-state. The yamabushi offered the body itself as the instrument of return---physical shock, sensory overwhelm, the authority of flesh over spirit.

Soken listened. He watched Ensho through the open door of the zendo, still sitting, still breathing, still seeing whatever he saw. The practitioners had come from different directions, carrying different maps. But the territory they described was the same. The eye that tries to see itself. The hearing that hears itself hearing. The seeing that sees itself and finds no ground.

They had all seen this before. Or heard of it. Or lost someone to it.

The oldest fear. The one that preceded every formal teaching, every lineage, every tradition: that the human mind, pushed far enough, could achieve a perception it could not survive.


\section*{III. The Aftermath}

They sat with him through the night.

Soken had sent the other monks away---to the kitchen, to the storage buildings, to their cells. Not because there was danger in proximity. Because what was happening in the zendo was private in a way that the young monks could not be expected to understand, and Soken would not make a spectacle of his student's departure.

Departure. He was already using the word in his mind, though Ensho's body remained present, breathing, warm to the touch when Soken pressed his fingers to the back of Ensho's hand. The hand did not respond. Not a flinch, not a withdrawal, not the subtle muscular acknowledgment that someone has been touched. The body continued on its own---heart, lungs, the slow work of digestion---but the person who lived inside it was attending to something else. Something that required all available attention and then more than all.

The four of them sat in a rough circle around Ensho. Lamplight threw long shadows across the cedar walls. Outside, the wind had dropped, and the silence was so complete that Soken could hear the snow settling on the temple roof---a sound like nothing, like the weight of air gaining substance.

``In my lineage,'' Wuzhun said, ``we do not interfere. We house them. We feed them. We observe. What he sees may be true. We have no right to end it.''

``Truth that consumes the one who sees it is not liberation,'' Dorje Wangchuk said, through the translation chain. ``It is another kind of trap. In the bardo, the mind encounters visions. Some are true. The instruction is the same: do not cling to them. Do not reject them. Observe and pass through. He has stopped passing through. He is clinging to what he sees.''

``He is not clinging,'' Soken said. ``What he sees is clinging to him.''

Silence. The distinction mattered. A mind that clings can be taught to release. A perception that holds the mind---that will not release because seeing it produces more seeing, and the more seeing produces more still---that is a different problem.

Kaido, the yamabushi, spoke from the edge of the circle. ``I have seen this in the mountains. Not during meditation. During the cold-water practice. A man sits under a waterfall in winter. His body becomes so cold that his mind separates from it. Usually the cold brings him back---the body screams loud enough that the mind returns. But sometimes the mind finds something during the separation. Something it cannot release. The man stays in the water. The body dies. The mind---'' He shrugged. ``Who knows what the mind does? It goes where we cannot follow.''

``Your student is not dying,'' Dorje Wangchuk observed.

``No,'' Soken said. ``He is not dying. He is not suffering. That is what frightens me.''

The admission cost him. In the Rinzai tradition, the teacher does not confess fear. The teacher strikes with the keisaku, roars the koan, demands the student demonstrate understanding with immediate, unrehearsed action. The teacher is the stone the student breaks against, and the stone does not admit to cracking. But Soken was too old and the situation too far beyond his training for pretense.

``What frightens me is that what he sees may be real. What frightens me is that I recognize it. Not from seeing it myself---I did not go that far. From recognizing the shape of it. The accounts my teacher gave me. His teacher before him. The descriptions that were always kept vague, deliberately, because---''

He stopped. Because he had been about to describe the thing that could not be described, and even in this company, even among practitioners who had seen similar territory, the instinct to protect was stronger than the need to explain.

``Because the description is the thing,'' Wuzhun finished. ``Yes. In our lineage as well.''

\begin{center}
*\quad*\quad*
\end{center}

Dorje Wangchuk tried his method first. He sat before Ensho and chanted in a low voice---Tibetan syllables, rhythmic, precise, the bardo instructions meant to guide a mind through the between-state. The words were not magic. They were a map spoken aloud, describing the landscape the lost mind was traversing, offering landmarks, suggesting a direction toward return.

Ensho's breathing changed. Slowed further. As if the chanting were not guiding him back but drawing him deeper---as if the map, perceived at his expanded depth, was revealing territory beyond what the map-maker had intended.

Dorje Wangchuk stopped. His face was troubled.

``He is hearing the instructions,'' the Tibetan said. ``But he is hearing more than the instructions. He is hearing the structure beneath the instructions. The pattern that makes the instructions meaningful. He is---'' The translation chain faltered. Dorje Wangchuk spoke rapidly in Tibetan, and by the time the meaning arrived in Japanese, it was compressed to a single sentence: ``He hears the hearing.''

The yamabushi's method was simpler. Kaido brought water from the mountain stream---ice-cold, carried in a wooden bucket. He poured it over Ensho's head and shoulders. The shock was immediate and physical: Ensho gasped, his body convulsing, muscles contracting against the cold. A reflex. The body remembering what the mind had forgotten. For one moment, Ensho was entirely present in his body---soaked, shaking, cold.

Then the shaking stopped. The breathing slowed. The eyes, which had closed against the water, opened again. The same gaze. The same attention directed at the thing with no location. The body had shouted and the mind had heard it and decided the body was not the important thing.

Soken watched. He thought of his teacher Genro, who had described the progression in careful, coded words: \emph{First they sit deeper than others. Then they hear what others do not hear. Then the hearing turns inward. Then the inward hearing hears itself. Then there is no more ``then,'' because ``then'' requires time, and time has become something else for them.}

Time has become something else. Soken had always taken this as metaphor. Looking at Ensho now, he was not certain it was.

\begin{center}
*\quad*\quad*
\end{center}

He sent the others out. Not because their help was unwelcome. Because the decision, whatever it was, was his. Ensho was his student. The debt ran in one direction: teacher to student, the oldest obligation in the lineage, older than any institution or network or tradition. You take a student. You guide the student toward awakening. If the awakening destroys the student, that is your failure. Not the student's. Not the tradition's. Yours.

The zendo was empty except for the two of them. The lamps had burned low. The cold had deepened. Soken could see his own breath, but he could not see Ensho's---the breathing had become so shallow, so slow, that no visible vapor formed.

Soken picked up the keisaku. The familiar weight. He had carried this stick for twenty years, had struck ten thousand blows across the shoulders of monks who needed waking, correcting, sharpening. The stick was an extension of the teaching. The teaching was: \emph{come back}. Come back to the body. Come back to the breath. Come back to the koan. Come back to the small, bounded, finite reality of one mind in one body in one moment.

He struck. Hard. Harder than he had ever struck anyone. The sound filled the zendo---a crack like a branch breaking under the weight of snow.

Ensho did not move.

Soken set the keisaku down. He sat before Ensho, knee to knee. He took Ensho's hands. They were cold but alive---the pulse was there, slow, steady, counting out a rhythm that belonged to whatever time Ensho was keeping.

``Ensho.'' He squeezed the hands. Hard enough that the bones compressed, hard enough to bruise. ``The mountain is cold tonight. The stream is frozen at the edges. Your robe is wet from the water Kaido poured. You will become ill if you do not change it. Ensho. These are small things. These are the things that a person tends to.''

Nothing. The eyes looking past him, through him, at the thing.

``Come back,'' Soken said. The words were not a command. They were an offering. Come back to this. To the cold room, the wet robe, the old man holding your hands. Come back to the bounded, the limited, the incomplete. Come back to the place where the eye does not see itself and the mercy of that blindness makes everything else possible.

A change. Barely. A shift in the quality of Ensho's gaze, as if the thing he was looking at had, for one moment, become slightly less compelling than the thing in front of him---the human face, the familiar room, the cold.

His eyes focused. On Soken. On the specific, particular reality of one old monk in a dark zendo.

Three heartbeats. Soken counted them by the pulse in the hands he held.

``Roshi,'' Ensho said. His voice was very far away, but it was his voice, and it was speaking to Soken, and for those three heartbeats he was present.

``Where does the seeing stop?'' Ensho asked.

The question was not rhetorical. It was genuine. He was asking because he did not know. He had followed the seeing---the seeing of seeing, the hearing of hearing, the awareness of awareness of awareness---and he had not found a bottom. He had not found a place where consciousness rested on something that was not itself consciousness. No floor. No ground. No place where the turning stopped.

Soken opened his mouth to answer. He had no answer. The tradition had no answer. The question was the oldest in the lineage and no teacher had ever resolved it, because resolving it required standing outside consciousness to observe its foundation, and there was no outside.

The three heartbeats ended. Ensho's eyes shifted. The gaze moved past Soken, past the walls of the zendo, past the mountain and the forest and the river and the thin autumn sky. To the thing. To the seeing that would not stop seeing itself.

He did not come back.

\begin{center}
*\quad*\quad*
\end{center}

They cared for him. There was no discussion about this; it was simply what was done. Wuzhun, who had done this before, showed the younger monks how to bring rice gruel to Ensho's lips, how to tip the bowl so the liquid entered his mouth and his body's reflexes---still intact, still operating beneath whatever was occupying his attention---would swallow. They changed his wet robe for a dry one. They placed a heavier blanket over his shoulders when the cold deepened.

The other monks adapted. Within a week, Ensho's motionless form in the back row of the zendo had become a feature of the room, like the altar or the cedar pillars. The youngest monk, the boy who had whispered his concern weeks ago, now swept around Ensho with the same unselfconscious care he used when sweeping around the furniture. Some of the monks, Soken noticed, sat straighter during zazen when they were near Ensho. As if his presence---his depth, his terrible achievement---raised the seriousness of the room.

Soken wrote letters. Not coded, this time. The first danger---the danger of description triggering perception---had passed. What remained was factual: a practitioner at the mountain temple had achieved a state from which he would not return. The temple would care for him. If other temples had experienced similar cases, the mountain temple would welcome correspondence.

Three replies arrived over the following months. From a temple in Echizen, a brief account of a nun, two generations ago, who had sat in zazen for forty days without eating or speaking. She died. The account was unsure whether her death was the failure or the success. From a monastery in the mountains above Kyoto, a longer account of a monk who had entered a similar state and been retrieved---pulled back by his teacher's intervention, though the account did not specify what the intervention was. The monk had lived another twenty years but had never sat zazen again. He gardened. He cooked. He swept. He was kind to the younger monks. He never spoke of what he had seen, and his eyes, the account noted, were ``ordinary'' afterward---as if whatever had expanded in them had been deliberately contracted, and the contraction was permanent.

From a Tibetan monastery, reached through Dorje Wangchuk's connections, a single line in Tibetan that the chain of translators rendered as: \emph{We know. We are sorry. It does not stop.}

\begin{center}
*\quad*\quad*
\end{center}

The visitors departed. Wuzhun returned to Kamakura, then to China. Dorje Wangchuk continued south with his trading party. Kaido went back to his mountains.

Soken remained. In the mornings, before the other monks woke, he sat in the zendo with Ensho. Not across from him---beside him. Two old men on cushions in the dark, one seeing and one choosing not to see. The keisaku lay across Soken's knees. He did not use it. He could not say whether his failure was in striking too late, or not striking hard enough, or in believing that striking was ever the right response to a student who had gone where the stick could not follow.

His own teacher's words circled in his mind like water in a basin: \emph{Some students the teacher pulls back. Some students the teacher pushes forward. The tragedy is that sometimes pulling back and pushing forward lead to the same place.}

He understood now. Ensho had not needed pulling or pushing. Ensho had needed thirty years of patience and a koan that would not resolve and the accumulated weight of ten thousand sittings, each one a fraction deeper, each one stretching the membrane between perception and the thing perceived until the membrane was no longer there. The direction---forward or back---was irrelevant. The membrane was gone. There was nothing to pull him back to.

Soken would pass this understanding to his own students. Carefully. In words that described the shape without revealing the substance. Fingers pointing at the moon. Warnings delivered in coded language that a prepared mind could decode and an unprepared mind would hear as ordinary teaching. The accounts would travel from teacher to student, generation to generation, each transmission losing some detail and preserving what mattered: this is possible. Watch for it. The eye cannot see itself. This is not a limitation. This is a mercy.

Whether the mercy was kindness or cowardice---whether the bounded, limited, incomplete perception of ordinary human consciousness was a gift to be grateful for or a prison to escape from---Soken did not know. He sat with the question the way Ensho had sat with his koan. He did not expect an answer. He did not abandon it.

Early winter. Snow on the temple roofs. The stream running thin and clear beneath a shelf of ice. In the zendo, Ensho breathed. His eyes tracked something the other monks could not see. The youngest monk brought him rice gruel each morning and each evening, lifting the bowl to lips that accepted without acknowledgment. The temple bells rang the hours. The seasons would turn. Ensho would sit. Soken would watch. The question would remain.

The moon was there. No finger had ever touched it. No finger ever would.
